\section{Further remarks}
    \label{sec:further-remarks}

It is not known whether the ratio of $M(N)$ and $N^2(\log N)^{-\delta} (\log \log N)^{-3/2}$ tends to a limit as $N \rightarrow \infty$, and it remains an open problem to find a true asymptotic formula for $M(N)$.
Computational experiments were conducted, among others, by Brent et al. \cite{algorithmsforthemtp}, but the slow growth of the function $(\log N)^\delta (\log \log N)^{3/2}$ makes it challenging to extract numerical evidence for whether an asymptotic formula exists.

Apart from the classical multiplication table problem, several generalisations have been considered in the literature. 
Notably, Koukoulopoulos \cite{localizedfactorizationsofintegers} adapted Ford's arguments to determine the order of magnitude of a $(k + 1)$-dimensional multiplication table, that is, the number of integers that can be written as a product of $k + 1$ integers $\le N$.
The analysis is based on a higher-dimensional analogue of $H(x, y, 2y)$.

Another possible generalisation is the \defn{restricted} multiplication table problem:
Given an infinite subset $\mathcal A \In \mN$, how many distinct elements of $\mathcal A$ can be written as a product of two integers $\le N$?
This corresponds to studying the quantity $H(x, y, z, \mathcal A)$, the number of integers $n \in \mathcal A \cap[1, x]$ that have a divisor in $(y, z]$.
In the special case where $\mathcal A = \set {p + b}{p \text{ prime}}$ for some $b \ne 0$, the order of magnitude was determined by Ford \cite{thedistributionofintegers} and Koukoulopoulos \cite{divisorsofshiftedprimes}.
